\clearpage
\section{Gaussian cellular automata}\label{sec:ccrautomata}

\noindent\motiv{\reve{Chapter~\ref{sec:gaussianqca} studied rules acting on the phase space itself; this chapter
lets them act instead on the observables. A symplectic rule $L$ defines an automorphism $\alpha_L$ of
the CCR algebra by $\alpha_L(W(u))=W(Lu)$, and this automorphism is local, commutes with the lattice
translations, and composes, so that it is a quantum cellular automaton in the sense made precise
below. Conversely, any automorphism sending Weyl operators to multiples of Weyl operators arises from
a symplectic linear bijection of $V$ together with a character, and the bijection is a rule
whenever the automorphism also commutes with the translations. The leapfrog again serves as the
running example, and at the end of the chapter it is identified with one time step of the lattice
Klein--Gordon field.}}

\medskip
\noindent\motiv{\reve{Throughout this chapter $G$ is an abelian group, the cell carries $n$ modes, and $V$ denotes
the lattice phase space of Definition~\ref{def:qca-space} over $\R$, equipped with the symplectic
form $\sigma=\tilde\Omega$. The algebra $\CCR(V,\sigma)$ is the C*-algebra generated by the Weyl
operators $W(u)$ on $\ell^2(V)$, which act by
\[(W(u)f)(v)=e^{\ii\sigma(u,v)/2}f(v+u),\]
and consequently satisfy
\[W(u)W(w)=e^{-\ii\sigma(u,w)/2}\,W(u+w),\qquad W(u)^*=W(-u).\]} \reve{Here $\ell^2(V)$ is the space of
square-summable functions on the set $V$, with the point masses $\delta_v$ as orthonormal basis; it
serves only to construct the algebra, and it is not the Fock space of the oscillators, since
$t\mapsto W(tu)$ is not continuous on it and it carries no Gaussian state
(Remark~\ref{rem:gs-existence}). The algebra, and not this representation, is what the states and
channels of later chapters act on.}}

\subsection{The lattice CCR net}

\noindent\motiv{\reve{Local observables are attached to regions of the lattice, understood as subsets
$O\subseteq G$: observables belonging to disjoint regions should commute, and the observables of a
region belong as well to every larger region containing it. Since the lattice carries no time
direction, there is no time-slice condition beyond this isotony.}}

\noindent\motiv{\reve{Every linear bijection of $V$ that preserves $\sigma$ relabels the Weyl operators, and both
the translations and the rules of the previous chapter are bijections of this kind.}}

\begin{lemma}[Relabelling]\label{lem:relabel}
Let $P\colon V\to V$ be a linear bijection with \reve{$\sigma(Pu,Pw)=\sigma(u,w)$}. The unitary
$(U_P\psi)(v)=\psi(P^{-1}v)$ on $\ell^2(V)$ satisfies $U_PW(u)U_P^*=W(Pu)$, so $\alpha_P=U_P\,\cdot\,U_P^*$ is a
$*$-automorphism of $\CCR(V,\sigma)$ with $\alpha_P(W(u))=W(Pu)$, and $\alpha_{PP'}=\alpha_P\alpha_{P'}$.
\end{lemma}

\begin{proof}
$U_P$ is unitary because $P$ is a bijection, and $(U_P^*\psi)(v)=\psi(Pv)$. Hence
\[(U_PW(u)U_P^*\psi)(v)=e^{\ii\sigma(u,P^{-1}v)/2}\,\psi(v+Pu)=e^{\ii\sigma(Pu,v)/2}\,\psi(v+Pu)=(W(Pu)\psi)(v).\]
Conjugation by $U_P$ is a $*$-automorphism of the bounded operators, isometric, and maps the Weyl
operators onto the Weyl operators; so it maps $\CCR(V,\sigma)$, the closure of their span, onto
itself. Finally $U_{PP'}=U_PU_{P'}$. Two automorphisms that agree on the Weyl operators agree on
their span, and so on its closure; so $\alpha_P$ is the only automorphism with
$\alpha_P(W(u))=W(Pu)$.
\end{proof}

\begin{definition}[The lattice CCR net]\label{def:ccr-net}
For $O\subseteq G$, $\cA(O)$ is the closed $*$-subalgebra of $\CCR(V,\sigma)$ generated by the $W(u)$
with $\supp u\subseteq O$. The translation by $a\in G$ acts by $\alpha_a=\alpha_{\kappa_{a*}}$, so
$\alpha_a(W(u))=W(\kappa_{a*}u)$.
\end{definition}

\begin{theorem}\label{thm:ccr-net}
The net $O\mapsto\cA(O)$ is isotone, generates $\CCR(V,\sigma)$, and satisfies causality: elements of
$\cA(O_1)$ and $\cA(O_2)$ commute when $O_1\cap O_2=\emptyset$. The map $a\mapsto\alpha_a$ is a group
homomorphism into the automorphisms of $\CCR(V,\sigma)$, and $\alpha_a\cA(O)=\cA(O+a)$.
\end{theorem}

\begin{proof}
\reve{Isotony holds because $\cA(O)$ has fewer generators than $\cA(O')$ whenever $O\subseteq O'$,
while $\cA(G)=\CCR(V,\sigma)$ by definition. For causality, if $\supp u\cap\supp w=\emptyset$ then
every term of $\sigma(u,w)=\sum_xu(x)^{\mathsf T}\Omega w(x)$ vanishes, so that $W(u)$ and $W(w)$ commute;
since the commutant of the generators of $\cA(O_2)$ is a closed $*$-algebra, being closed under $*$ as
a set, and contains the generators of $\cA(O_1)$, it contains $\cA(O_1)$ itself, and by the same
argument the commutant of $\cA(O_1)$ contains $\cA(O_2)$. For the last statement, translations
preserve $\sigma$, so $\alpha_a$ exists by Lemma~\ref{lem:relabel}, and $\alpha_{a+b}=\alpha_a\alpha_b$
because $\kappa_{(a+b)*}=\kappa_{a*}\kappa_{b*}$; since $\supp\kappa_{a*}u=\supp u+a$, $\alpha_a$
carries the generators of $\cA(O)$ onto those of $\cA(O+a)$, and being an isometric $*$-automorphism
it therefore carries $\cA(O)$ onto $\cA(O+a)$.}
\end{proof}

\subsection{The automorphism of a rule}

\noindent\motiv{\reve{Since a symplectic rule $L$ is a bijection of $V$ preserving $\sigma$
(Theorem~\ref{thm:qca-inverse}), Lemma~\ref{lem:relabel} already furnishes its automorphism
$\alpha_L$.}}

\begin{proposition}\label{prop:alpha}
For symplectic rules $L$ and $K$:
\begin{enumerate}[label=(\alph*)]
\item $\alpha_L(W(u))=W(Lu)$;
\item $\alpha_{LK}=\alpha_L\alpha_K$ and $\alpha_1=\mathrm{id}$;
\item $\alpha_L\alpha_a=\alpha_a\alpha_L$ for every $a\in G$;
\item $\alpha_L^{-1}=\alpha_{L^{-1}}$, where $L^{-1}$ has symbol $\Omega^{-1}\hat L^{\star\mathsf T}\Omega$.
\end{enumerate}
\end{proposition}

\begin{proof}
\begin{enumerate}[label=(\alph*)]
\item \reve{Lemma~\ref{lem:relabel}.}
\item \reve{Lemma~\ref{lem:relabel}.}
\item $L$ commutes with $\kappa_{a*}$ (Theorem~\ref{thm:qca-symbol}(a)), so both sides send $W(u)$ to
$W(L\kappa_{a*}u)$.
\item By (b) and Theorem~\ref{thm:qca-inverse}.
\end{enumerate}
\end{proof}

\begin{theorem}[Locality]\label{thm:alpha-local}
Let $N$ be the range of the symbol of $L$. Then $\alpha_L\cA(O)\subseteq\cA(O+N)$ for every $O\subseteq G$.
\end{theorem}

\begin{proof}
If $\supp u\subseteq O$, then $\supp Lu\subseteq O+N$ (Corollary~\ref{cor:qca-range}), so $\alpha_L$ maps
the generators of $\cA(O)$ into $\cA(O+N)$. A $*$-homomorphism maps the $*$-algebra generated by a set
onto the $*$-algebra generated by its image; a $*$-automorphism of a C*-algebra is isometric, so it
maps closures into closures, and $\cA(O+N)$ is closed.
\end{proof}

\subsection{Quantum cellular automata}

\begin{definition}[Reversible quantum cellular automaton]\label{def:rqca}
Let $\cA$ be a net over the subsets of $G$ in a C*-algebra $A$, with a group homomorphism
$a\mapsto\alpha_a$ into the automorphisms of $A$ such that $\alpha_a\cA(O)=\cA(O+a)$. A \emph{reversible quantum cellular automaton} is a
$*$-automorphism $T$ of $A$ together with a finite set $N\subseteq G$, its \emph{neighbourhood}, such
that $T\alpha_a=\alpha_aT$ for all $a$ and $T\cA(O)\subseteq\cA(O+N)$ for all $O$.
\end{definition}

\noindent\motiv{\reve{The locality condition $T\cA(O)\subseteq\cA(O+N)$ requires that a
single step change the observables of a region only within the neighbourhood of that region.}}

\begin{proposition}[Composition]\label{prop:rqca-comp}
If $(T_1,N_1)$ and $(T_2,N_2)$ are reversible quantum cellular automata on the same net, so is
$(T_2T_1,N_1+N_2)$.
\end{proposition}

\begin{proof}
$T_2T_1$ commutes with every $\alpha_a$, and $T_2T_1\cA(O)\subseteq T_2\cA(O+N_1)\subseteq\cA(O+N_1+N_2)$.
\end{proof}

\begin{theorem}\label{thm:rule-qca}
For a symplectic rule $L$ with range $N$, $(\alpha_L,N)$ is a reversible quantum cellular automaton
on the lattice CCR net.
\end{theorem}

\begin{proof}
Proposition~\ref{prop:alpha}(c) and Theorem~\ref{thm:alpha-local}; $N$ is finite because the entries of
the symbol are finitely supported. \reve{The inverse $\alpha_{L^{-1}}$ is likewise a reversible quantum cellular automaton, with
neighbourhood $-N$, since the entries of $\Omega^{-1}\hat L^{\star\mathsf T}\Omega$ are, up to sign,
exactly those of $\hat L^{\star\mathsf T}$.}
\end{proof}

\begin{example}[The leapfrog]\label{ex:ccr-leapfrog}
The leapfrog $F$ of Example~\ref{ex:qca-leapfrog-rule} and its inverse give reversible quantum
cellular automata $\alpha_F$ and $\alpha_{F^{-1}}$ with neighbourhood $\{-1,0,1\}$, and
$\alpha_F^{-1}=\alpha_{F^{-1}}$. In the Cauchy data $u(x)=(d_0(x),d_1(x))$,
$\alpha_F(W(u))=W(Fu)$ with $(Fu)(x)=\big(d_1(x),\,-d_0(x)+d_1(x+1)+d_1(x-1)-m^2d_1(x)\big)$.
\end{example}

\begin{example}[Shifts]\label{ex:ccr-shift}
For $a\in G$, the symbol $z^a\cdot1$ gives the reversible quantum cellular automaton with
neighbourhood $\{a\}$ whose automorphism is the translation $\alpha_a$ itself.
\end{example}

\begin{example}[On-site rules]\label{ex:ccr-onsite}
For $G=\Z$ and a real $2n\times2n$ matrix $S$ with $S\Omega S^{\mathsf T}=\Omega$, the constant symbol $S$ gives a
reversible quantum cellular automaton with neighbourhood $\{0\}$: every site evolves separately by
$S$. For one mode this covers the rotations $R_\varphi=\begin{psmallmatrix}\cos\varphi&-\sin\varphi\\
\sin\varphi&\cos\varphi\end{psmallmatrix}$, the free evolution of uncoupled oscillators, and the
squeezings $\operatorname{diag}(e^r,e^{-r})$, which iterated squeeze every site further.
\end{example}

\begin{example}[Rotation with symmetric coupling]\label{ex:ccr-rotation}
The rule of Example~\ref{ex:qca-rotation} gives a reversible quantum cellular automaton with
neighbourhood $\{-1,0,1\}$ for all $\varphi$ and $\lambda$. \reve{A neighbourhood need not be the smallest possible one, as is seen already at $\lambda=0$,
where the range of the symbol reduces to $\{0\}$.}
\end{example}

\subsection{The local rule}

\noindent\motiv{\reve{A rule is entirely determined by its action on the degrees of freedom of a single site. For
$i\in\iota$, let $u_i\in V$ denote the vector taking the value $e_i$ at the site $0$ and vanishing
elsewhere, so that the operators $W(tu_i)$, $t\in\R$, are precisely the Weyl operators of the site
$0$.}}

\begin{proposition}\label{prop:local-rule}
\begin{enumerate}[label=(\alph*)]
\item A \reve{translation-invariant linear map} is determined by the vectors $Lu_i$, $i\in\iota$.
\item A \reve{translation-invariant linear map} $L$ is symplectic if and only if
\[\sigma(Lu_i,\kappa_{x*}Lu_j)=\delta_{x0}\,\Omega_{ij}\qquad\text{for all }i,j\in\iota,\ x\in G.\]
\end{enumerate}
\end{proposition}

\begin{proof}
\begin{enumerate}[label=(\alph*)]
\item $Lu_i$ is the $i$-th column of the symbol \reve{(the vector whose value at $x\in G$ is $\hat L_{ji}$, read as a function on $G$), and the
symbol in turn determines $L$ by Theorem~\ref{thm:qca-symbol}.}
\item $\sigma(Lu_i,\kappa_{x*}Lu_j)$ is the coefficient of $(\hat L^{\star\mathsf T}\Omega\hat L)_{ij}$ at
$-x$, as in the proof of Theorem~\ref{thm:qca-criterion}; the condition says that this symbol is the
constant $\Omega$.
\end{enumerate}
\end{proof}

\noindent\motiv{\reve{In terms of the automaton, part (b) says that the images $\alpha_L(W(tu_i))=W(tLu_i)$ of the
one-site Weyl operators, together with their translates, satisfy the same commutation relations as
the one-site Weyl operators themselves.}}

\subsection{Automorphisms of Weyl type}

\noindent\motiv{\reve{The automorphisms of $\CCR(V,\sigma)$ that send every Weyl operator to a multiple of a Weyl
operator are characterized by means of a norm estimate. Let
$\delta_0\in\ell^2(V)$ be the unit vector at $0\in V$; the functional $\varphi(a)=(a\delta_0)(0)$
satisfies $|\varphi(a)|\le\|a\|$, and $\varphi(W(u))=1$ if $u=0$ and $0$ otherwise, since
$(W(u)\delta_0)(0)=\delta_0(u)$.}}

\begin{lemma}[Separation]\label{lem:separation}
If $y$ lies in the closed linear span of $\{W(v):v\neq u\}$, then $\|W(u)-y\|\ge1$. In particular,
$\mu W(u)=\nu W(w)$ with $\nu\neq0$ implies $u=w$ and $\mu=\nu$.
\end{lemma}

\begin{proof}
For $v\neq u$, $W(u)^*W(v)$ is a multiple of $W(v-u)$, so $\varphi(W(u)^*y')=0$ for $y'$ in the span and,
by continuity, in its closure, while $\varphi(W(u)^*W(u))=1$. Hence
$1=|\varphi(W(u)^*(W(u)-y))|\le\|W(u)-y\|$. For the second statement apply $\varphi(W(w)^*\,\cdot\,)$ to
both sides.
\end{proof}

\begin{theorem}\label{thm:weyl-type}
Let $T$ be a $*$-automorphism of $\CCR(V,\sigma)$, $Q\colon V\to V$ a map and $c\colon V\to\C$ with
$T(W(u))=c(u)W(Q u)$ for all $u$. Then $Q$ is a linear bijection with
$\sigma(Q u,Q w)=\sigma(u,w)$, and $c$ is a character: $c(u+w)=c(u)c(w)$, $|c(u)|=1$.
\end{theorem}

\begin{proof}
$c(u)\neq0$ since $T$ is injective, and $|c(u)|=1$ since $T$ is isometric and Weyl operators have
norm $1$.

\emph{Additivity.} Applying $T$ to $W(u)W(w)=e^{-\ii\sigma(u,w)/2}W(u+w)$ gives
\[c(u)c(w)\,e^{-\ii\sigma(Q u,Q w)/2}\,W(Q u+Q w)=e^{-\ii\sigma(u,w)/2}\,c(u+w)\,W(Q(u+w)),\]
so $Q(u+w)=Q u+Q w$ and the coefficients agree, by Lemma~\ref{lem:separation}.

\emph{Symplecticity.} Exchanging $u$ and $w$ and dividing gives $e^{-\ii\sigma(Q u,Q w)}=
e^{-\ii\sigma(u,w)}$, so $\sigma(u,w)-\sigma(Q u,Q w)=2\pi k(u,w)$ with $k(u,w)\in\Z$. By additivity,
$m\,Q(u/m)=Q u$ for $m\in\mathbb N$, hence $k(u,w)=m\,k(u/m,w)$; so $k(u,w)$ is divisible by every
$m$, and $k=0$.

\emph{Surjectivity.} If $v'\notin Q V$, then $T$ maps the span of the Weyl operators into the span
of $\{W(v):v\neq v'\}$. The Weyl operators span a dense subspace of $\CCR(V,\sigma)$ (their span is
already a $*$-algebra), and $T$ is continuous, so $W(v')=T(T^{-1}W(v'))$ lies in the closure of that
span, contradicting Lemma~\ref{lem:separation}.

\emph{Linearity.} For $t\in\R$, $\sigma(Q(tu)-tQ u,Q w)=\sigma(tu,w)-t\sigma(u,w)=0$ for every $w$;
since $Q$ is onto and $\sigma$ is nondegenerate (if $\sigma(u,w)=0$ for all $w$, pairing with the
one-site vectors $z^xe_j$ gives $\Omega^{\mathsf T} u(x)=0$ for every $x$, so $u=0$), $Q(tu)=tQ u$. If $Q u=0$ then
$\sigma(u,w)=0$ for all $w$, so $u=0$.

\emph{Character.} With symplecticity, the coefficient identity of the first step reads
$c(u)c(w)=c(u+w)$.
\end{proof}

\noindent\motiv{\reve{Nontrivial phases $c$ do occur in Theorem~\ref{thm:weyl-type}. For any character $c$ with $|c|=1$, the unitary
$(U\psi)(v)=\overline{c(v)}\,\psi(v)$ satisfies $UW(u)U^*=c(u)W(u)$, so that $\alpha_L$ followed by
conjugation with $U$ is an automorphism of the form of the theorem, with $Q=L$ and phase $c\circ L$.
When $T$ also commutes with the translations, comparing $T\alpha_a(W(u))$ with $\alpha_aT(W(u))$ shows
that $Q$ commutes with $\kappa_{a*}$, so that $Q$ is itself a rule and $c\circ\kappa_{a*}=c$. Such
translation-invariant characters do exist, for instance
$c(u)=\exp\big(\ii\sum_x\theta^{\mathsf T}u(x)\big)$ with $\theta\in\R^{2n}$, and conjugation with the
corresponding $U$ is a reversible quantum cellular automaton with neighbourhood $\{0\}$. The
translation-covariant automata of Weyl type are therefore exactly the automorphisms $\alpha_L$
followed by such a conjugation.}}

\subsection{The leapfrog and the direction of time}

\noindent\motiv{\reve{The Cauchy data $(d_0,d_1)$ of the lattice Klein--Gordon field form exactly the one-mode
lattice phase space over $\Z$, through $u(x)=(d_0(x),d_1(x))$, and their conserved form
$\sum_x\big(d_0(x)d_1'(x)-d_1(x)d_0'(x)\big)$ is $\sigma$. The leapfrog step $F$ of
Example~\ref{ex:qca-leapfrog-rule} advances the data from the slices $t,t+1$ to $t+1,t+2$, and it
remains to determine which automaton represents time evolution itself.}}

\begin{proposition}\label{prop:time-direction}
Let $\Lambda_t$ send a solution $\varphi$ of the lattice Klein--Gordon equation \reve{with finitely
supported Cauchy data} to its Cauchy data
$(\varphi(t,\cdot),\varphi(t+1,\cdot))$, and let $\tau$ translate solutions forward in time by one step,
$(\tau\varphi)(t,x)=\varphi(t-1,x)$. Then $\Lambda_t\tau=F^{-1}\Lambda_t$.
\end{proposition}

\begin{proof}
$\Lambda_t\tau\varphi=(\varphi(t-1,\cdot),\varphi(t,\cdot))=\Lambda_{t-1}\varphi$, and $\Lambda_t=F\Lambda_{t-1}$ by the
construction of $F$ in Example~\ref{ex:qca-leapfrog-rule}. Hence $\Lambda_{t-1}=F^{-1}\Lambda_t$.
\end{proof}

\noindent\motiv{\reve{Translating a solution forward in time shifts its Cauchy data at a fixed time backwards, and
on observables the consequence of this shift depends on how $W(u)$ is to be read. If $W(u)$ is
thought of as $e^{\ii\sigma(u,\Phi_t)}$, with $\Phi_t$ the Cauchy data of the quantum field at time
$t$ (a heuristic reading here, since $\ell^2(V)$ carries no field operators and $t\mapsto W(tu)$ is
not continuous), then the field moves as $\Phi_{t+1}=F\Phi_t$, so that
$\sigma(u,\Phi_{t+1})=\sigma(F^{-1}u,\Phi_t)$; under this reading the forward time step is therefore
the automaton $\alpha_{F^{-1}}$ of Example~\ref{ex:ccr-leapfrog}, rather than $\alpha_F$.}}

\subsection{Notes and references}\label{sec:ccr-notes}

\begin{enumerate}
\item \reve{\emph{Definition.} Definition~\ref{def:rqca} corresponds to Definition~4.1 of~\cite[p.~73]{krueger}
for reversible automata, with the quasi-local algebra there replaced by a net; the condition
$0\in N$ imposed in the thesis is not needed here. Irreversible automata, that is, quantum channels,
are treated \reve{in Chapter~\ref{sec:channels}}.}
\item \reve{\emph{Composition and local rule.} Proposition~\ref{prop:rqca-comp} corresponds to
Corollary~4.3 of the thesis. Proposition~\ref{prop:local-rule}(b) is the Gaussian analogue of
Lemma~4.2(ii) rather than a special case of it: that lemma, like the results
of~\cite{schumacher-werner}, concerns quasi-local algebras built from finite-dimensional cells,
whereas each cell here carries the Weyl algebra of $\R^{2n}$.}
\item \reve{\emph{Weyl-type automorphisms.} After (4.7), \cite[p.~78]{krueger} states that any $\Gamma$
for which $T(W(\xi))=W(\Gamma\xi)$ defines a homomorphism ``would have to be linear \dots\ and
fulfill $\sigma(\Gamma\xi,\Gamma\eta)=\sigma(\xi,\eta)$''. Theorem~\ref{thm:weyl-type} proves this claim
and additionally records the character $c$, which the thesis omits. Lemma~\ref{lem:relabel}, together
with the uniqueness of the CCR algebra for nondegenerate $\sigma$ (established in the proof of
Theorem~\ref{thm:weyl-type}), of which $\CCR(V,\sigma)$ is a concrete realization, is due to
Manuceau~\cite{manuceau}, Slawny~\cite{slawny}, and Manuceau, Sirugue, Testard and
Verbeure~\cite{mstv}; see~\cite[\S5.2]{bratteli-robinson}.}
\item \reve{\emph{Examples.} The thesis lists the free evolution, the left- and right-shifter, continued
squeezing, and symplectic rotations as examples of Gaussian automata (p.~70);
Example~\ref{ex:ccr-rotation} corresponds to its system (4.18).}
\item \reve{\emph{Direction of time.} In the convention $T(W(\xi))=W(\Gamma\xi)$ used in the thesis, the
rule governing the forward time step of the lattice Klein--Gordon field is $F^{-1}$.}
\end{enumerate}
